\section{Теоретические основы машинно-обучаемых потенциалов}
\label{sec:theory}

\subsection{Постановка задачи регрессии энергий и сил}

Атомную конфигурацию будем записывать как набор
\[
    X = \{(\mathbf r_i, z_i)\}_{i=1}^{N},
\]
где \(\mathbf r_i \in \mathbb R^3\) -- координаты атома, \(z_i\) -- тип атома, \(N\) -- число атомов в конфигурации. Машинно-обучаемый межатомный потенциал строит приближение к потенциальной энергии \(E^{\mathrm{ref}}(X)\) и силам \(\mathbf F^{\mathrm{ref}}_{i}(X)\), взятым из эталонного расчёта или подготовленного датасета.

В локальных потенциалах полная энергия представляется как сумма атомных вкладов:
\[
    E_\theta(X) = \sum_{i=1}^{N} V_\theta(\mathcal N_i),
\]
где \(\theta\) -- параметры модели, \(\mathcal N_i\) -- локальное окружение атома \(i\) в пределах радиуса отсечения. Силы вычисляются как производные энергии по координатам:
\[
    \mathbf F_{i,\theta}(X) =
    -\frac{\partial E_\theta(X)}{\partial \mathbf r_i}.
\]

Обучение далее рассматривается как минимизация функционала ошибки на обучающей выборке \(\{X_s\}_{s=1}^{S}\):
\[
\mathcal L(\theta)
=
w_E \sum_{s=1}^{S}
\left(
E_\theta(X_s) - E_s^{\mathrm{ref}}
\right)^2
+
w_F \sum_{s=1}^{S}
\sum_{i=1}^{N_s}
\left\|
\mathbf F_{i,\theta}(X_s)
-
\mathbf F_{i,s}^{\mathrm{ref}}
\right\|^2 .
\]
В MLIP-4 значение loss и метрики ошибок вычисляются объектом \texttt{LossFunction}; в результатах дополнительно используются RMSE по полной энергии, энергии на атом (EPA RMSE) и силам.

\subsection{Moment Tensor Potential}

Moment Tensor Potential (MTP) относится к локальным межатомным потенциалам: атомная энергия в нём раскладывается по инвариантным базисным функциям локального окружения \cite{shapeev2016mtp,novikov2021mlip}. Для атома \(i\) вводится относительный радиус-вектор
\[
    \mathbf r_{ij} = \mathbf r_j - \mathbf r_i .
\]
Моментный тензор определяется формулой
\[
M_{\mu,\nu}(\mathcal N_i)
=
\sum_{j \in \mathcal N_i}
f_\mu
\left(
|\mathbf r_{ij}|, z_i, z_j
\right)
\underbrace{
\mathbf r_{ij} \otimes \cdots \otimes \mathbf r_{ij}
}_{\nu \text{ раз}},
\]
где \(\mu\) индексирует радиальную часть, \(\nu\) задаёт тензорный порядок, а \(z_i,z_j\) задают типы центрального и соседнего атомов.

Радиальная часть имеет вид
\[
f_\mu(r, z_i, z_j)
=
\sum_{k=1}^{K}
c_{\mu k z_i z_j}
\varphi_k(r),
\]
где \(\varphi_k(r)\) -- радиальные базисные функции, а \(c_{\mu k z_i z_j}\) -- обучаемые коэффициенты.

Базисные инварианты строятся как полные свёртки моментных тензоров:
\[
B_\alpha(\mathcal N_i)
=
\operatorname{contr}
\left(
M_{\mu_1,\nu_1},
M_{\mu_2,\nu_2},
\dots,
M_{\mu_m,\nu_m}
\right),
\]
где \(\operatorname{contr}\) обозначает полную свёртку тензорных индексов, дающую скалярный инвариант. Атомная энергия MTP задаётся линейной комбинацией таких инвариантов:
\[
V_\theta(\mathcal N_i)
=
\sum_{\alpha}
\xi_\alpha
B_\alpha(\mathcal N_i),
\qquad
E_\theta(X)
=
\sum_{i=1}^{N}
V_\theta(\mathcal N_i).
\]

Тензор \(M_{\mu,\nu}\) одновременно несёт радиальную и угловую информацию о локальном окружении. При \(\nu=0\) он даёт скалярный признак, при \(\nu=1\) -- векторный, при \(\nu=2\) -- тензор второго порядка. Последующие свёртки переводят эти объекты в скалярные функции, инвариантные относительно поворотов конфигурации. Параметр \texttt{level} регулирует сложность MTP: при его увеличении в модель попадают более сложные произведения и свёртки моментных тензоров.

\subsection{Equivariant Tensor Network}

Второй класс моделей в работе -- эквивариантный тензорно-сетевой потенциал ETN \cite{orus2014tnintro,stoudenmire2016tnml,hodapp2025etn}. Как и MTP, он использует локальное разложение энергии по атомным окружениям:
\begin{equation}
    E_\theta(X)
    =
    \sum_{i=1}^{N}
    V_\theta(\mathcal N_i),
\end{equation}
где \(X=\{(\mathbf r_i,z_i)\}_{i=1}^{N}\) --- атомная конфигурация, \(\mathbf r_i\) --- координаты атома, \(z_i\) --- тип атома, \(\mathcal N_i\) --- локальное окружение атома \(i\) в пределах заданного радиуса отсечения, а \(V_\theta(\mathcal N_i)\) --- атомный вклад в энергию.

Главное отличие ETN от MTP связано со способом построения признаков локального окружения. В MTP заранее задаётся набор моментных тензоров и инвариантных базисных функций, а атомная энергия записывается как их линейная комбинация. В ETN локальное окружение сначала кодируется набором эквивариантных признаков, а затем эти признаки сворачиваются в скалярную атомную энергию с помощью тензорной сети.

Для пары атомов \(i,j\) вводится относительный радиус-вектор
\begin{equation}
    \mathbf r_{ij}
    =
    \mathbf r_j - \mathbf r_i,
    \qquad
    r_{ij}
    =
    \|\mathbf r_{ij}\|,
    \qquad
    \widehat{\mathbf r}_{ij}
    =
    \frac{\mathbf r_{ij}}{r_{ij}}.
\end{equation}
Радиальная зависимость описывается набором радиальных базисных функций
\begin{equation}
    \varphi_k(r_{ij}),
    \qquad
    k=1,\ldots,K.
\end{equation}
В используемых экспериментах применялся радиальный базис \texttt{RadialBasisCinf} с параметрами \texttt{size=8}, \texttt{min\_dist=1.9}, \texttt{cutoff=5.0}. Химические типы атомов учитываются через зависимость коэффициентов модели от \(z_i\) и \(z_j\).

Эквивариантные признаки можно схематически записать в виде
\begin{equation}
    h^{(\ell)}_{i,c,m}
    =
    \sum_{j\in \mathcal N_i}
    \sum_{k}
    a^{(\ell)}_{c k}(z_i,z_j)
    \varphi_k(r_{ij})
    Y_{\ell m}(\widehat{\mathbf r}_{ij}),
\end{equation}
где \(\ell\) --- угловой порядок признака, \(m=-\ell,\ldots,\ell\) --- индекс компоненты в неприводимом представлении вращений, \(c\) --- номер канала, \(Y_{\ell m}\) --- сферические гармоники, а \(a^{(\ell)}_{c k}(z_i,z_j)\) --- обучаемые коэффициенты смешивания радиальных и химических признаков. Эта формула не фиксирует все детали реализации ETN, но отражает основной принцип: радиальная часть зависит от расстояний, а угловая часть задаёт поведение признаков при поворотах.

Ключевое свойство таких признаков состоит в эквивариантности. Если вся конфигурация поворачивается ортогональной матрицей \(Q\), то признаки порядка \(\ell\) преобразуются согласованно с представлением группы вращений:
\begin{equation}
    h^{(\ell)}_{i,c}(QX)
    =
    D^{(\ell)}(Q)
    h^{(\ell)}_{i,c}(X),
\end{equation}
где \(D^{(\ell)}(Q)\) --- матрица представления поворота \(Q\) для углового порядка \(\ell\). При \(\ell=0\) признак является скаляром и не меняется при повороте. При \(\ell=1\) признак преобразуется как вектор, а при больших \(\ell\) --- как тензор более высокого углового порядка.

Атомная энергия должна быть инвариантной величиной. Поэтому в ETN эквивариантные признаки сворачиваются так, чтобы итоговый результат имел нулевой угловой порядок:
\begin{equation}
    V_\theta(\mathcal N_i)
    =
    \operatorname{TN}_\theta
    \left(
        \{h^{(\ell)}_{i,c}\}_{\ell,c}
    \right)
    \in \mathbb R.
\end{equation}
Здесь \(\operatorname{TN}_\theta\) обозначает тензорную сеть с обучаемыми параметрами. Промежуточные признаки могут быть векторными или тензорными, но итоговая атомная энергия остаётся скаляром. Поэтому полная энергия удовлетворяет условию инвариантности:
\begin{equation}
    E_\theta(QX)
    =
    E_\theta(X).
\end{equation}

Силы вычисляются как производные энергии по координатам:
\begin{equation}
    \mathbf F_{i,\theta}(X)
    =
    -
    \frac{\partial E_\theta(X)}
    {\partial \mathbf r_i}.
\end{equation}
Из инвариантности энергии следует корректное преобразование сил при повороте конфигурации:
\begin{equation}
    \mathbf F_{i,\theta}(QX)
    =
    Q\mathbf F_{i,\theta}(X).
\end{equation}

В практической реализации ETN задаётся несколькими группами гиперпараметров. Параметр \texttt{etn\_shape} определяет ширину набора каналов для разных уровней или порядков признаков. В базовых экспериментах использовалось
\begin{equation}
    \texttt{etn\_shape}
    =
    [[1],[6],[3]].
\end{equation}
Увеличение \texttt{etn\_shape} повышает выразительность модели, но одновременно увеличивает число параметров и усложняет задачу оптимизации. Параметр \texttt{ett\_ranks} задаёт ранги тензорной сети:
\begin{equation}
    \texttt{ett\_ranks}
    =
    [[1],[1]].
\end{equation}
Малые ранги дают более компактную модель, а большие ранги увеличивают число степеней свободы. Поэтому \texttt{etn\_shape} и \texttt{ett\_ranks} влияют не только на аппроксимационную способность ETN, но и на то, насколько сложной становится задача оптимизации.

В итоге ETN можно рассматривать как модель, где физические симметрии встроены непосредственно в архитектуру. В отличие от MTP, где набор инвариантных базисных функций задан явно, ETN использует тензорно-сетевое смешивание эквивариантных признаков. Такая конструкция делает модель более гибкой, но одновременно усложняет ландшафт функции потерь. Поэтому результаты, полученные для MTP, нельзя автоматически переносить на ETN; именно это проверяется в экспериментальной части.
